\documentclass[11pt,a4paper]{article}
\usepackage[margin=1in]{geometry}
\usepackage{amsmath, amssymb}
\usepackage{xcolor}

\definecolor{primary}{RGB}{0, 105, 148}
\definecolor{secondary}{RGB}{10, 150, 120}

\setlength{\parindent}{0pt}
\setlength{\parskip}{0.5em}

\begin{document}

\begin{center}
    \Large\textbf{\textcolor{primary}{Introduction to Markov Chains}}
\end{center}

\vspace{0.5cm}

\section*{\textcolor{secondary}{Stochastic Process:}}
A stochastic process is simply a collection of random variables, usually indexed by time: $\{X_t, t \in T\}$\\
where, $t$ is time (discrete or continuous)\\
$X_t$ is the state of the process at time $t$

\# The set of all possible values that $X_t$ can take is called the \textbf{state space}.

\section*{\textcolor{secondary}{Markov Property (Memorylessness):}}
A sequence of random variables is a Markov Chain if it possesses the Markov Property.

Mathematically:
\[ P(X_{n+1} = j \mid X_n = i, X_{n-1} = i_{n-1}, \dots, X_0 = i_0) = P(X_{n+1} = j \mid X_n = i) \]

i.e. The probability of moving to state $j$ tomorrow ($X_{n+1}$), given the entire history of the system from day zero, is exactly the same as the probability of moving to state $j$ given only where we are today ($X_n = i$).

\section*{\textcolor{secondary}{One-Step Transition Probabilities ($P_{ij}$):}}
The probability of transitioning from state $i$ to state $j$ in exactly one time step is $P_{ij}$.
\[ P_{ij} = P(X_{n+1} = j \mid X_n = i) \]

These $P_{ij}$'s should satisfy two rules:
\begin{enumerate}
    \item Non-negativity: $P_{ij} \geq 0$
    \item Row sums to 1: $\sum_{j=0}^{\infty} P_{ij} = 1 \quad \forall i$
\end{enumerate}

\section*{\textcolor{secondary}{Transition Probability Matrix ($P$):}}
All these individual transition probabilities when packaged into a single square matrix ($P$)
\[ P = \begin{bmatrix} P_{00} & P_{01} & P_{02} & \dots \\ P_{10} & P_{11} & \dots \\ \vdots & \vdots \\ P_{i0} & P_{i1} & \dots \\ \vdots & \vdots \end{bmatrix} \]

\vspace{0.5cm}

\section*{\textcolor{primary}{Chapman - Kolmogorov Equations:}}
\textbf{\textcolor{secondary}{Core Problem:}} We know $P_{ij}$ which is probability of going from $i$ to $j$ in exactly 1 step.

But what if we want to know the probability of going from $i$ to $j$ in exactly $n$ steps?\\
We denote this as $P_{ij}^{(n)}$.

\textbf{\textcolor{secondary}{Equation:}} \\
To figure out the probability of going from $i$ to $j$ in $n+m$ steps, you have to account for every possible intermediate state $k$.

Formally:
\[ P_{ij}^{(n+m)} = \sum_{k=0}^{\infty} P_{ik}^{(n)} P_{kj}^{(m)} \]

If we let $P^{(n)}$ denote the matrix of n-step transition probabilities $P_{ij}^{(n)}$, then
\[ P^{(n+m)} = P^{(n)} \cdot P^{(m)} \]
$\hookrightarrow$ matrix multiplication

\section*{\textcolor{primary}{Examples}}
From Introduction to Probability Models by Sheldon Ross:

\textbf{Example 4.10.} An urn always contains 2 balls. Ball colors are red and blue. At each
stage a ball is randomly chosen and then replaced by a new ball, which with probability 0.8 is the same color, and with probability 0.2 is the opposite color, as the ball it
replaces. If initially both balls are red, find the probability that the fifth ball selected is
red.

\textbf{Solution:}
\begin{minipage}[t]{0.5\textwidth}
\textbf{states:} \\
state 0 $\rightarrow$ 0 Red 2 Blue \\
state 1 $\rightarrow$ 1 Red 1 Blue \\
state 2 $\rightarrow$ 2 Red 0 Blue
\end{minipage}
\begin{minipage}[t]{0.45\textwidth}
Let $X_n$ be the no. of red balls in urn after $n$th selection and subsequent replacement.
\end{minipage}

\vspace{0.5cm}
\textbf{Transition Matrix:}
\[ P = \begin{bmatrix} 0.8 & 0.2 & 0 \\ 0.1 & 0.8 & 0.1 \\ 0 & 0.2 & 0.8 \end{bmatrix} \]

P(fifth selection is red)
\begin{align*}
    &= \sum_{i=0}^2 P(\text{fifth selection is red} \mid X_4 = i) P(X_4 = i \mid X_0 = 2) \\
    &= (0) P_{2,0}^{(4)} + (0.5) P_{2,1}^{(4)} + (1) P_{2,2}^{(4)} \\
    &= 0.7048 \quad (\text{after carrying matrix multiplications})
\end{align*}

\newpage

\textbf{Example 4.11.} Suppose that balls are successively distributed among 8 urns, with
each ball being equally likely to be put in any of these urns. What is the probability
that there will be exactly 3 nonempty urns after 9 balls have been distributed?

\textbf{Solution:} \\
$\rightarrow$ $X_n$ be the no of nonempty urns after $n$ balls have been distributed. \\
states: $0$ to $8$ \\
$P_{i,i} = i/8$ \\
$P_{i,i+1} = 1 - P_{i,i}$

We want to find $P_{0,3}^{(9)}$.
\[ P_{0,3}^{(9)} = P_{0,1}^{(1)} \cdot P_{1,3}^{(8)} = 1 \cdot P_{1,3}^{(8)} = P_{1,3}^{(8)} \]

To build a transition matrix, we have states $1, 2, 3, 4$ (in fact, $\geq 4$ nonempty urns).
\[ P = \begin{bmatrix} 1/8 & 7/8 & 0 & 0 \\ 0 & 2/8 & 6/8 & 0 \\ 0 & 0 & 3/8 & 5/8 \\ 0 & 0 & 0 & 1 \end{bmatrix} \]
we need $P_{1,3}^{(8)}$.

After Matrix multiplications
\[ P_{1,3}^{(8)} = 0.00756 \]

\vspace{3em}

\textbf{Example 4.12.} In a sequence of independent flips of a fair coin, let $N$ denote the number of flips until there is a run of three consecutive heads. Find \\
(a) $P(N \leq 8)$ and \\
(b) $P(N = 8)$.

\textbf{Solution:} \\
$\rightarrow$ Markov chain, as we are currently on a run of $i$ consecutive heads. \\
states: $0, 1, 2, 3$ ($3 \rightarrow$ 3 consecutive heads have already occurred)

\[ P = \begin{bmatrix} 1/2 & 1/2 & 0 & 0 \\ 1/2 & 0 & 1/2 & 0 \\ 1/2 & 0 & 0 & 1/2 \\ 0 & 0 & 0 & 1 \end{bmatrix} \]

We want: \\
a) $P_{0,3}^{(8)}$ \\
b) $P_{0,2}^{(7)} P_{2,3}^{(1)}$

Calculate $P^7$ and $P^8$ to get the answer.

\end{document}